\documentclass[12pt]{article}
\usepackage[T1]{fontenc}
\usepackage[utf8]{inputenc}
\usepackage{lmodern}
\usepackage{textcomp}
\usepackage{enumitem}
\usepackage{geometry}
\geometry{margin=1in}
\begin{document}
\begin{center}
Answer Key - Philosophy - Undergraduate Level Assessment 12 \
\small (Answers are ordered exactly as in the question paper.)
\end{center}

\section*{Multiple Choice}
\begin{enumerate}[label=\arabic*.]
\item \textbf{Answer:} A
\\ \textit{Options:} A) peace.; B) virtue.; C) Godliness.; D) wisdom.; E) happiness.
\\ \textit{Note:} Aquinas identifies peace as the ultimate perfection of operation because it represents the harmonious state of being in which all desires and actions align with one's rational nature and the good, ultimately leading to fulfillment and well-being.
\item \textbf{Answer:} B
\\ \textit{Options:} A) a priori; B) Complex proposition; C) Extension; D) Appeal to ignorance
\\ \textit{Note:} The correct answer is "Complex proposition" because this fallacy occurs when an argument attempts to validate an entire proposition by providing evidence or support for only one component, thereby misleadingly asserting that the whole proposition is substantiated.
\item \textbf{Answer:} D
\\ \textit{Options:} A) The parents are prioritizing their own needs over the child's.; B) We can't tell what the motives of the parents are.; C) The parents have ulterior motives.; D) The parents have good motives.; E) The parents are incapable of making decisions for the child.
\\ \textit{Note:} Pence argues that those who object to somatic cell nuclear transfer (SCNT) assume that the parents involved in the process are driven by altruistic motives, prioritizing the child's well-being over any potential ethical concerns.
\item \textbf{Answer:} A
\\ \textit{Options:} A) strong; B) reliant on circular reasoning; C) too complicated; D) overly simplistic; E) biased in favor of Cleanthes' conclusion
\\ \textit{Note:} Philo asserts that Cleanthes' analogy effectively illustrates the existence of a divine creator through its compelling and coherent comparison between the universe and human-made objects, thus deeming it a strong argument in favor of theism.
\item \textbf{Answer:} C
\\ \textit{Options:} A) Laudatory personality; B) Red herring; C) Reprehensible personality; D) Circular reasoning; E) Straw man fallacy
\\ \textit{Note:} The "reprehensible personality" fallacy occurs when one dismisses a person's positive actions solely based on their negative traits, illustrating a flawed reasoning that equates character with the ability to perform good deeds.
\end{enumerate}

\section*{True / False}
\begin{enumerate}[label=\arabic*.]
\setcounter{enumi}{5}
\item \textbf{Answer:} True
\\ \textit{Options:} A) True; B) False
\\ \textit{Note:} "Ad lazarum" is indeed a type of fallacy where an argument is based on misleading evidence, leading one to draw false conclusions, as it relies on the erroneous assumption that a lack of evidence or a flawed sign means the opposite is true.
\item \textbf{Answer:} True
\\ \textit{Options:} A) True; B) False
\\ \textit{Note:} Defenders of meat-eating claim that human rationality allows for a moral distinction between human suffering and animal suffering, suggesting that humans can weigh consequences and make ethical decisions, which justifies the consumption of meat.
\item \textbf{Answer:} False
\\ \textit{Options:} A) True; B) False
\\ \textit{Note:} The statement is false because, while Taiwan values Confucian principles, the largest network of Confucian shrines is actually found in mainland China, where Confucianism has deeper historical roots and more extensive institutional support.
\item \textbf{Answer:} False
\\ \textit{Options:} A) True; B) False
\\ \textit{Note:} The statement is false because, despite some examples of strong female figures in Hindu mythology and history, traditional interpretations and societal practices often depict women in subordinate roles, emphasizing their duties in the domestic sphere rather than positions of empowerment and leadership.
\item \textbf{Answer:} False
\\ \textit{Options:} A) True; B) False
\\ \textit{Note:} The statement is false because Altman argues that while causal harm may justify speech codes, expressive harm does not provide a sufficient basis for restricting speech.
\end{enumerate}

\section*{Long Form Response}
\begin{enumerate}[label=\arabic*.]
\setcounter{enumi}{10}
\item \textbf{Answer:} To construct a truth table for the argument represented by the premises \textasciitilde{}S ∨ T, \textasciitilde{}S · U, and \textasciitilde{}T ∨ U, we can analyze the truth values of S, T, and U systematically. The truth table reveals combinations of truth values for S, T, and U, ultimately showing that the premises do not guarantee a true conclusion in all scenarios. Specifically, the argument is invalid; a counterexample occurs when U is true, while both S and T are false. In this case, the premises \textasciitilde{}S ∨ T and \textasciitilde{}S · U are true, but the conclusion derived from the premises does not hold, demonstrating that the argument does not maintain logical validity across all interpretations.
\\ \textit{Note:} correct because the constructed truth table demonstrates that there exists a scenario where the premises are true but the conclusion is false, indicating that the argument is invalid, with the specific counterexample provided illustrating this discrepancy.
\item \textbf{Answer:} The fallacy of extension occurs when an argument improperly generalizes a claim or attribute from one context to another, misapplying the original intention behind the subject. This fallacy undermines the validity of a critique, particularly when the subject was never intended to fulfill the criticized purpose, as it misrepresents the argument's original context and goals. For instance, if a philosophical theory is critiqued for failing to address practical application when it was designed solely as a theoretical framework, the critique becomes irrelevant and misleading. Consequently, this fallacy not only distorts the argument but also detracts from constructive discourse, as it leads to misunderstandings and misinterpretations of the subject's actual aims and relevance. Thus, recognizing and avoiding the fallacy of extension is essential for maintaining the integrity of philosophical discussions and critiques.
\\ \textit{Note:} correct because it accurately identifies how the fallacy of extension misapplies a subject's original intent, leading to an invalid critique that fails to consider the appropriate context and purpose of the argument being evaluated.
\end{enumerate}

\vfill
\noindent\textit{End of Answer Key}
\end{document}